\section{introduction}
Parton distributions characterize the structure of nucleons in terms of partons in high-energy scattering processes. They are an essential ingredient in making physical predictions for hadron-hadron or lepton-hadron collision experiments. Although much effort has been devoted to extracting parton distributions from various experimental data~\cite{Alekhin:2012ig,Gao:2013xoa,Radescu:2010zz,CooperSarkar:2011aa,Martin:2009iq,Ball:2012cx}, the computation of parton distributions from the underlying theory of strong interactions, quantum chromodynamics (QCD), is a difficult task, due to their non-perturbative nature. One may wonder whether it is possible to evaluate parton distributions from lattice QCD, which is so far the only reliable framework for non-perturbative phenomena in QCD. In the field-theoretic language, the parton distribution is given in terms of non-local light-cone correlators, which can be viewed as the density of quarks and gluons in the infinite momentum limit (before ultraviolet (UV) cut-offs are imposed) or light-front
correlations of partons at finite nucleon momentum~\cite{Collins:2011zzd}. Such correlations are time-dependent, and thus cannot be readily computed on the lattice. Past attempts have been mainly focused on the evaluation of local moments of the distributions. However, the evaluation of high moments of parton distributions on the lattice becomes considerably difficult for technical reasons~\cite{negele}.

Recently, a direct approach to compute parton physics on a Euclidean lattice
has been proposed by one of the present authors. In this approach one computes, instead of the light-cone distribution, a related quantity, which may be called quasi-distribution~\cite{footnote}. In the case of unpolarized
quark density, the quasi-distribution is~\cite{Ji:2013dva}
\begin{eqnarray}\label{qUnpolDef}
   \tilde q(x,\mu^2, P^z) &=& \int \frac{dz}{4\pi} e^{izk^z}  \langle P|\overline{\psi}(z)
   \gamma^z \\
   && \times \exp\left(-ig\int^{z}_0 dz' A^z(z') \right)\psi(0) |P\rangle \ ,  \nonumber
\end{eqnarray}
where $x= k^z/P^z$. The above quantity is time-independent and non-local, and can be simulated on a lattice
for any $P^z\ll 1/a$, where $a$ is the lattice spacing. However, the result is not the light-cone distribution extracted
from the experimental data, $q(x, \mu^2)$. To recover the latter from the former,
one needs to find a matching condition of the type (for the non-singlet case)
\begin{equation}
   \tilde q_{\rm NS}(x, \mu^2, P^z) = \int dy\, Z(x/y, \mu/P^z) q_{\rm NS}(y, \mu^2) + {\cal O} ((M/P^z)^2) \ ,
\end{equation}
for large $P^z$, where the matching factor $Z$ is entirely perturbative. The above relation can also
be viewed as a factorization theorem: All the soft divergences are cancelled on both sides, and all
the collinear divergences in $\tilde q(x, \mu^2, P^z)$ are the same as those in the light-cone
distributions. One has yet to prove that the above relation holds to all orders in perturbation
theory. However, as a first step, we will test this in one-loop calculations in this paper. Of course the choice of quasi-distributions is not unique, one can define more than one possible quasi-distributions which have similar properties
as $\tilde q(x, \mu^2, P^z)$. Here we will focus on the simplest type suitable for lattice QCD calculations.

In this paper, we consider the non-singlet quark distribution, and prove the above factorization at one-loop level. Throughout the process we obtain the matching factor $Z$ between the time-independent quasi-distribution and light-cone distribution. The result will be useful for
recovering the light-cone distribution from lattice QCD simulations of the quasi-distribution to the leading
logarithmic accuracy. Power suppressed corrections of the type $(1/P^z)^{n}$ with $n\ge 1$ are ignored and will be considered
in separate publications.

The rest of this paper is organized as follows. In Sec. 2, we present the result of one-loop calculation for unpolarized quasi-quark distribution. Based on this result, we propose in Sec. 3 a factorization theorem valid to one-loop level and extract the matching factor between quasi distribution and light-cone distribution.
In Sec. 4, the results for quark helicity and transversity distributions are given. Finally we conclude in Sec. 5.
